% GENERATED FILE. Edit canonical lessons, then run npm run generate:books.
% canonical-source-hash: fnv1a64:f9d9dbeb148e9c75
% canonical-lessons: TE-C391-pramukhyata, TE-C391-pradhanyam, TE-C391-prathamika, TE-C391-nokkiceppu, TE-C391-prastavincu

\chapter{What Matters Most}
\label{ch:te-what-matters-most}
\hlchaptermodality{391}

\begin{chapteropening}
\textbf{By the end of this chapter:} \emph{I can say what is important, basic and worth stressing.}

Complete the last lesson of chapter 391: I can say what is important, basic and worth stressing.
\end{chapteropening}

\section[\texorpdfstring{\te{ప్రాముఖ్యత} (prāmukhyata)}{ప్రాముఖ్యత (prāmukhyata)}]{\te{ప్రాముఖ్యత} (prāmukhyata) — importance; significance}

\label{lesson:TE-C391-pramukhyata}



\begin{quote}
Before the new one: say the Telugu for a process, then the Telugu for to define.
\end{quote}

\subsection*{You'll want to know: \te{ప్రాముఖ్యత}}
\textbf{\te{ప్రాముఖ్యత}} — \emph{prāmukhyata} — \textquotedblleft{}importance; significance\textquotedblright{}.

\begin{quote}
\te{ఆరోగ్యానికి} \te{మనం} \te{ఎక్కువ} \te{ప్రాముఖ్యత} \te{ఇవ్వాలి}.
\end{quote}

\begin{quote}
\emph{ārōgyāniki manaṁ ekkuva prāmukhyata ivvāli.}
\end{quote}

\begin{quote}
We should give more importance to health.
\end{quote}

\te{ప్రాముఖ్యత} \te{ఇచ్చు}, to give importance; the thing takes -\te{కి}.

\begin{cousinweb}
\textbf{\te{ప్రాముఖ్యత}} (\emph{prāmukhyata}) is built on Sanskrit \emph{pramukha}, \textquotedblleft{}facing; first, foremost, chief\textquotedblright{} (Monier-Williams), with the abstract ending -\te{త}.
\end{cousinweb}

\begin{grammarlens}[title={What you've built}]
One more word for giving reasons and opinions.
\end{grammarlens}

\subsection*{Your turn}
\begin{itemize}
\raggedright
  \item \emph{Say it:} \emph{prāmukhyata}
  \item \emph{Say it:} \emph{prāmukhyata}, once more
  \item \emph{Say it:} \emph{ārōgyāniki manaṁ ekkuva prāmukhyata ivvāli.} — We should give more importance to health.
  \item \emph{Recall:} say the Telugu for a process, then the Telugu for to define, then say \emph{prāmukhyata} again
\end{itemize}

\begin{tcolorbox}[breakable,colback=teal!4,colframe=teal!35!black,title={Before you move on}]
What does \te{ప్రాముఖ్యత} mean? (\textquotedblleft{}Importance; significance\textquotedblright{}.) Say it once more.
\end{tcolorbox}
\section[\texorpdfstring{\te{ప్రాధాన్యం} (prādhānyaṁ)}{ప్రాధాన్యం (prādhānyaṁ)}]{\te{ప్రాధాన్యం} (prādhānyaṁ) — priority; precedence}

\label{lesson:TE-C391-pradhanyam}



\begin{quote}
Before the new one: say the Telugu for to define, then the Telugu for importance.
\end{quote}

\subsection*{You'll want to know: \te{ప్రాధాన్యం}}
\textbf{\te{ప్రాధాన్యం}} — \emph{prādhānyaṁ} — \textquotedblleft{}priority; precedence\textquotedblright{}.

\begin{quote}
\te{అమ్మ} \te{పిల్లల} \te{చదువుకు} \te{ప్రాధాన్యం} \te{ఇస్తుంది}.
\end{quote}

\begin{quote}
\emph{amma pillala caduvuku prādhānyaṁ istundi.}
\end{quote}

\begin{quote}
Mother gives priority to the children’s studies.
\end{quote}

Close to \te{ప్రాముఖ్యత}, but it puts one thing ahead of the others.

\begin{cousinweb}
\textbf{\te{ప్రాధాన్యం}} (\emph{prādhānyaṁ}) is borrowed from Sanskrit \emph{prādhānya}, \textquotedblleft{}predominance, supremacy\textquotedblright{} (Monier-Williams), from \emph{pradhāna}, \textquotedblleft{}chief\textquotedblright{}; the Telugu adjective \te{ప్రధాన}, main, is the same word.
\end{cousinweb}

\begin{grammarlens}[title={What you've built}]
One more word for giving reasons and opinions.
\end{grammarlens}

\subsection*{Your turn}
\begin{itemize}
\raggedright
  \item \emph{Say it:} \emph{prādhānyaṁ}
  \item \emph{Say it:} \emph{prādhānyaṁ}, once more
  \item \emph{Say it:} \emph{amma pillala caduvuku prādhānyaṁ istundi.} — Mother gives priority to the children’s studies.
  \item \emph{Recall:} say the Telugu for to define, then the Telugu for importance, then say \emph{prādhānyaṁ} again
\end{itemize}

\begin{tcolorbox}[breakable,colback=teal!4,colframe=teal!35!black,title={Before you move on}]
What does \te{ప్రాధాన్యం} mean? (\textquotedblleft{}Priority; precedence\textquotedblright{}.) Say it once more.
\end{tcolorbox}
\section[\texorpdfstring{\te{ప్రాథమిక} (prāthamika)}{ప్రాథమిక (prāthamika)}]{\te{ప్రాథమిక} (prāthamika) — basic; primary}

\label{lesson:TE-C391-prathamika}



\begin{quote}
Before the new one: say the Telugu for importance, then the Telugu for priority.
\end{quote}

\subsection*{You'll want to know: \te{ప్రాథమిక}}
\textbf{\te{ప్రాథమిక}} — \emph{prāthamika} — \textquotedblleft{}basic; primary\textquotedblright{}.

\begin{quote}
\te{ప్రాథమిక} \te{పాఠశాల} \te{మా} \te{ఇంటి} \te{దగ్గర} \te{ఉంది}.
\end{quote}

\begin{quote}
\emph{prāthamika pāṭhaśāla mā iṇṭi daggara undi.}
\end{quote}

\begin{quote}
The primary school is near our house.
\end{quote}

It stands before a noun and does not change: \te{ప్రాథమిక} \te{పాఠశాల}, a primary school.

\begin{cousinweb}
\textbf{\te{ప్రాథమిక}} (\emph{prāthamika}) is borrowed from Sanskrit \emph{prāthamika}, \textquotedblleft{}belonging or relating to the first\textquotedblright{} (Monier-Williams), from \emph{prathama}, \textquotedblleft{}first\textquotedblright{}.
\end{cousinweb}

\begin{grammarlens}[title={What you've built}]
One more word for giving reasons and opinions.
\end{grammarlens}

\subsection*{Your turn}
\begin{itemize}
\raggedright
  \item \emph{Say it:} \emph{prāthamika}
  \item \emph{Say it:} \emph{prāthamika}, once more
  \item \emph{Say it:} \emph{prāthamika pāṭhaśāla mā iṇṭi daggara undi.} — The primary school is near our house.
  \item \emph{Recall:} say the Telugu for importance, then the Telugu for priority, then say \emph{prāthamika} again
\end{itemize}

\begin{tcolorbox}[breakable,colback=teal!4,colframe=teal!35!black,title={Before you move on}]
What does \te{ప్రాథమిక} mean? (\textquotedblleft{}Basic; primary\textquotedblright{}.) Say it once more.
\end{tcolorbox}
\section[\texorpdfstring{\te{నొక్కిచెప్పు} (nokkiceppu)}{నొక్కిచెప్పు (nokkiceppu)}]{\te{నొక్కిచెప్పు} (nokkiceppu) — to stress; to emphasise}

\label{lesson:TE-C391-nokkiceppu}



\begin{quote}
Before the new one: say the Telugu for priority, then the Telugu for basic.
\end{quote}

\subsection*{You'll want to know: \te{నొక్కిచెప్పు}}
\textbf{\te{నొక్కిచెప్పు}} — \emph{nokkiceppu} — \textquotedblleft{}to stress; to emphasise\textquotedblright{}.

\begin{quote}
\te{వైద్యుడు} \te{నిద్ర} \te{ప్రాముఖ్యతను} \te{నొక్కిచెప్పాడు}.
\end{quote}

\begin{quote}
\emph{vaidyuḍu nidra prāmukhyatanu nokkiceppāḍu.}
\end{quote}

\begin{quote}
The doctor stressed the importance of sleep.
\end{quote}

Past (I …): \te{నొక్కిచెప్పాను}, \emph{nokkiceppānu}.

\te{నొక్కి}, having pressed, + \te{చెప్పు}, to say: to say a thing while pressing on it.

\begin{cousinweb}
\textbf{\te{నొక్కిచెప్పు}} (\emph{nokkiceppu}) is native Dravidian, a compound of \emph{nokku}, \textquotedblleft{}to press, to squeeze\textquotedblright{} (DEDR 3780), and \emph{ceppu}, \textquotedblleft{}to say, to tell\textquotedblright{} (DEDR 1955).
\end{cousinweb}

\begin{grammarlens}[title={What you've built}]
One more word for giving reasons and opinions.
\end{grammarlens}

\subsection*{Your turn}
\begin{itemize}
\raggedright
  \item \emph{Say it:} \emph{nokkiceppu}
  \item \emph{Say it:} \emph{nokkiceppu}, once more
  \item \emph{Say it:} \emph{vaidyuḍu nidra prāmukhyatanu nokkiceppāḍu.} — The doctor stressed the importance of sleep.
  \item \emph{Recall:} say the Telugu for priority, then the Telugu for basic, then say \emph{nokkiceppu} again
\end{itemize}

\begin{tcolorbox}[breakable,colback=teal!4,colframe=teal!35!black,title={Before you move on}]
What does \te{నొక్కిచెప్పు} mean? (\textquotedblleft{}To stress; to emphasise\textquotedblright{}.) Say it once more.
\end{tcolorbox}
\section[\texorpdfstring{\te{ప్రస్తావించు} (prastāviñcu)}{ప్రస్తావించు (prastāviñcu)}]{\te{ప్రస్తావించు} (prastāviñcu) — to bring up; to refer to}

\label{lesson:TE-C391-prastavincu}



\begin{quote}
Before the new one: say the Telugu for basic, then the Telugu for to stress.
\end{quote}

\subsection*{You'll want to know: \te{ప్రస్తావించు}}
\textbf{\te{ప్రస్తావించు}} — \emph{prastāviñcu} — \textquotedblleft{}to bring up; to refer to\textquotedblright{}.

\begin{quote}
\te{సమావేశంలో} \te{అతను} \te{జీతాల} \te{విషయం} \te{ప్రస్తావించాడు}.
\end{quote}

\begin{quote}
\emph{samāvēśaṁlō atanu jītāla viṣayaṁ prastāviñcāḍu.}
\end{quote}

\begin{quote}
At the meeting he brought up the matter of salaries.
\end{quote}

Past (I …): \te{ప్రస్తావించాను}, \emph{prastāviñcānu}.

The noun is \te{ప్రస్తావన}, a mention.

\begin{cousinweb}
\textbf{\te{ప్రస్తావించు}} (\emph{prastāviñcu}) is built on Sanskrit \emph{prastāva}, \textquotedblleft{}introducing a topic, preliminary mention, allusion\textquotedblright{} (Monier-Williams), with the Telugu verb ending -\te{ఇంచు}.
\end{cousinweb}

\begin{grammarlens}[title={What you've built}]
One more word for giving reasons and opinions.
\end{grammarlens}

\subsection*{Your turn}
\begin{itemize}
\raggedright
  \item \emph{Say it:} \emph{prastāviñcu}
  \item \emph{Say it:} \emph{prastāviñcu}, once more
  \item \emph{Say it:} \emph{samāvēśaṁlō atanu jītāla viṣayaṁ prastāviñcāḍu.} — At the meeting he brought up the matter of salaries.
  \item \emph{Recall:} say the Telugu for basic, then the Telugu for to stress, then say \emph{prastāviñcu} again
\end{itemize}

\begin{tcolorbox}[breakable,colback=teal!4,colframe=teal!35!black,title={Before you move on}]
What does \te{ప్రస్తావించు} mean? (\textquotedblleft{}To bring up; to refer to\textquotedblright{}.) Say it once more.
\end{tcolorbox}
